% Ad Atticum 14.18 --- headnote
% ad-atticum-14-18, 9 May 44 BC, in Pompeiano

\textit{Cicero to Atticus, written at the Pompeian villa on 9 May
44 BC --- Perseus dateline \textit{Scr. in Pompeiano vii Id. Mai.
a. 710 (44)}. Atticus has been needling Cicero for praising
Dolabella's recent action (the dispersal of the Caesar-cult
gathering at the Forum column) too lavishly, and Cicero half
concedes the point: the praise was warm because Atticus's own
letters had set the temperature. But Dolabella is now estranged
from Atticus on the same financial grounds that have made him
Cicero's bitterest creditor --- the dowry of Tullia, owed since
the Kalends of January and still unpaid, even after Dolabella has
discharged his other enormous debts (through the convenient
agency of Faberius, Antony's secretary, and ``aid from Ops,''
that is, Caesar's confiscated treasury). Cicero permits himself
the pun \textit{opem ab Ope petierit} --- help sought from Ops,
goddess of plenty, by way of her temple where the funds were
stored --- and signals that he is not quite undone by sending a
barbed letter that Dolabella will struggle to face in person.}

\textit{The rest is briskly elliptical: business with Atticus's
agents (the Albian and Patulcian debts, Eros's failings,
Montanus, Servius's despair as he leaves for the East), and a
political coda. If Brutus will not come to the Senate on the
Kalends of June --- the meeting Antony is convoking --- Cicero
cannot see what he is to do in the Forum either. The Ides of
March, on the evidence now coming in, did not get the
Liberators very far; and Cicero is thinking, more and more, of
Greece. He plans to leave Pompeii on 10 May (\textit{vi Idus
Mai.}).}
